% ==============================================================================
% RESUMEN / ABSTRACT
% ==============================================================================
% Instrucciones de la plantilla:
% - Redactar en español e inglés (exactamente 1 párrafo para cada idioma).
% - Contenido sugerido:
%   * Objetivo principal de la práctica de laboratorio.
%   * Metodología empleada y herramientas/bibliotecas utilizadas.
%   * Principales resultados alcanzados y aprendizajes obtenidos.
% ==============================================================================

\begin{abstract}
En esta práctica de laboratorio implementamos el modelo de iluminación básico de Phong para representar superficies con apariencia mate y plástica mediante OpenGL 3.3. Para lograrlo, calculamos a nivel de shaders las componentes ambiental, difusa y especular, permitiendo que una fuente de luz se mueva dinámicamente en el plano XZ utilizando las teclas del teclado. Como resultado, los objetos de la escena tridimensional adquirieron diferentes aspectos de iluminación, colores y contraste, simulando la interacción de la luz sobre distintas superficies.
 
\vspace{0.3cm}
\noindent\textbf{Palabras clave:} Modelos de iluminación, Phong, shaders, componentes ambiental, difusa y especular, OpenGL.
\end{abstract}

\vspace{0.5cm}

\renewcommand{\abstractname}{Abstract}
\begin{abstract}
In this laboratory exercise, we implemented the basic Phong lighting model using OpenGL 3.3 to render surfaces with a matte, plastic-like appearance. To achieve this, we calculated the ambient, diffuse, and specular components within the shaders, allowing a light source to move dynamically across the XZ plane via keyboard input. As a result, objects within the 3D scene exhibited varied lighting, colors, and contrast, simulating the interaction of light with different types of surfaces.

\vspace{0.3cm}
\noindent\textbf{Keywords:} Illumination models, Phong, shaders, ambient, diffuse and specular components, OpenGL.
\end{abstract}

\vspace{0.6cm}
